\documentclass[titlepage,a4paper]{jsarticle}
\usepackage{../../../sty/import}% 各種パッケージインポート
\usepackage{../../../sty/title}% タイトルページの変更
\usepackage{../../../sty/answergrid}
%% タイトルページの変数
% レポートタイトル
\title{心理学特論第12回目出席票}
% 提出日
\expdate{\today}
% 科目名
\subject{心理学特論}
% 分野
\class{情報経営システム工学分野}
% 学年
\grade{M1}
% 学籍番号
\mynumber{24336488}
% 記述者
\author{本間三暉}

\begin{document}
% titleページ作成
\maketitle
\section*{keyword\#1}
ホランダーの方略
\section*{keyword\#2}
モスコビッチの方略
\section{「考えの壁」＝事前に想定するのと，実際にその局面で最適の行動ができるかどうかは別のものです。
自分の体験や他人の事例で「考えの壁」にぶち当たったと思えることがあれば，それをなるべくわかりやすく説明して下さい。}

\section{社会情勢の変化や技術革新，価値観の多様化によって「今日の正解が明日の失敗の種」となることもあり得ます。
従来は良しとされていたものが，次第に受け入れられなくなる，陳腐化する，または「失敗」に転じてしまったものを，現在のニュースや過去の事例から取り上げ，その概要をまとめて下さい。
（広義では，技術の発展に伴って廃れたものも含みます）}

\newpage
\section{心理学特論課題サイコロワーク9回答記入欄}
\AnswerGridFilled{
  {論理的誤差}
  {寛大化傾向}
  {気分一致効果}
  {ハインリッヒの法則}
  {デブリーフィング}
  {汎適応症候群}
  {ホーソン効果}
  {レジリエンス}
  {適応障害}
  {外傷後ストレス障害}
  {ペーシング}
  {閉ざされた質問}
  {開かれた質問}
  {傾聴}
  {パターナリズム}
  {コーチング}
  {メンター}
  {エンハンシング効果}
  {アンダーマイニング効果}
  {囚人のジレンマ}
}

\end{document}